\documentclass[11pt,a4paper]{article}

% Mise en page et typographie
\usepackage[utf8]{inputenc}
\usepackage[T1]{fontenc}
\usepackage[french]{babel}
\usepackage{lmodern}
\usepackage{geometry}
\usepackage{microtype}

% Notations scientifiques et extraits de code
\usepackage{amsmath}
\usepackage{siunitx}
\usepackage[version=4]{mhchem}
\usepackage{xcolor}
\usepackage{listings}
\usepackage{booktabs}
\usepackage{hyperref}

\geometry{margin=2.4cm}
\sisetup{locale=FR, per-mode=symbol}
\DeclareSIUnit\molecule{molécule}
\DeclareSIUnit\atm{atm}

\lstset{
  language=Python,
  basicstyle=\ttfamily\small,
  backgroundcolor=\color{black!4},
  frame=single,
  breaklines=true,
  showstringspaces=false
}

\hypersetup{
  colorlinks=true,
  linkcolor=blue!45!black,
  urlcolor=blue!55!black
}

\newcommand{\code}[1]{\texttt{#1}}

\title{Documentation de \code{1\_section\_efficace\_CO2\_15um.py}}
\author{Projet climat 2026}
\date{}

\begin{document}

\maketitle
\tableofcontents

\section{Objectif et résultat du programme}

Ce script est un prototype consacré à la bande d'absorption du dioxyde de
carbone, \ce{CO2}, autour de \SI{15}{\micro\metre}. Il utilise les raies
spectroscopiques de la base HITRAN pour calculer une \emph{section efficace
d'absorption} $\sigma$ en \si{\metre\squared\per\molecule}. Cette grandeur
représente la capacité d'une molécule à absorber le rayonnement à une
longueur d'onde donnée.

Le programme :

\begin{enumerate}
  \item initialise une base de données HAPI locale ;
  \item télécharge les raies de l'isotopologue principal du \ce{CO2} ;
  \item calcule la section efficace avec un profil de Voigt ;
  \item convertit le résultat de
        \si{\centi\metre\squared\per\molecule} en
        \si{\metre\squared\per\molecule} ;
  \item interpole la section efficace sur une grille de longueurs d'onde ;
  \item affiche le résultat entre 12 et \SI{18}{\micro\metre}.
\end{enumerate}

Le script affiche un graphique, mais ne l'enregistre pas dans un fichier.

\section{Installation et accès aux données}

L'environnement du projet est prévu pour Python~3.12. Depuis le dossier
contenant le script et le fichier \code{requirements.txt} :

\begin{lstlisting}[language=bash]
python3.12 -m venv .venv
source .venv/bin/activate
python -m pip install -r requirements.txt
\end{lstlisting}

Les trois dépendances externes sont :

\begin{itemize}
  \item \code{numpy}, pour construire la grille spectrale et interpoler les
        valeurs ;
  \item \code{matplotlib}, pour produire le graphique ;
  \item \code{hitran-api}, qui s'importe dans Python sous le nom \code{hapi},
        pour télécharger les raies HITRAN et calculer la section efficace.
\end{itemize}

\code{pathlib} n'apparaît pas dans \code{requirements.txt}, car ce module fait
partie de la bibliothèque standard de Python.

\paragraph{Connexion internet.}
L'appel à \code{fetch} est exécuté à chaque lancement du script. Dans sa
forme actuelle, le programme tente donc de contacter le serveur HITRAN à
chaque exécution, même si des fichiers \code{CO2.data} et \code{CO2.header}
existent déjà. Ces fichiers constituent bien une copie locale, mais le code ne
vérifie pas leur présence avant d'appeler \code{fetch}.

\section{Principe physique et conventions}

HITRAN décrit le spectre en fonction du nombre d'onde
$\widetilde{\nu}$, exprimé en \si{\per\centi\metre}. Le graphique est, lui,
tracé en fonction de la longueur d'onde $\lambda$. Pour une longueur d'onde
exprimée en mètres :

\begin{equation}
  \widetilde{\nu}=\frac{1}{100\lambda}.
\end{equation}

Le facteur 100 convertit les mètres en centimètres. Ainsi,

\begin{equation}
  \lambda=\SI{15}{\micro\metre}
  \quad\Longrightarrow\quad
  \widetilde{\nu}\simeq\SI{667}{\per\centi\metre}.
\end{equation}

Le profil de Voigt employé par HAPI combine l'élargissement Doppler, lié au
mouvement thermique des molécules, et l'élargissement collisionnel, lié à la
pression. La section calculée dépend donc notamment de la température et de la
pression fournies à HAPI.

\section{Imports du programme}

Le fichier importe uniquement les objets dont il a besoin :

\begin{lstlisting}
from pathlib import Path

import matplotlib.pyplot as plt
import numpy as np
from hapi import absorptionCoefficient_Voigt, db_begin, fetch
\end{lstlisting}

\subsection{\code{Path} depuis \code{pathlib}}

\code{Path} sert à construire un chemin indépendant du système
d'exploitation. Il est utilisé dans :

\begin{lstlisting}
DOSSIER_DONNEES = Path(__file__).resolve().parent / "donnees_hitran"
\end{lstlisting}

\code{\_\_file\_\_} désigne le fichier Python en cours d'exécution,
\code{resolve()} produit son chemin absolu, \code{parent} sélectionne son
dossier et l'opérateur \code{/} ajoute le sous-dossier
\code{donnees\_hitran}. L'import de \code{Path} est donc utile : il garantit
que le dossier de données est recherché à côté du script, quel que soit le
dossier depuis lequel la commande Python est lancée.

\subsection{\code{numpy} sous l'alias \code{np}}

Deux fonctions NumPy sont appelées :

\begin{itemize}
  \item \code{np.linspace(12e-6, 18e-6, 3000)} crée 3000 longueurs d'onde
        régulièrement espacées entre 12 et \SI{18}{\micro\metre} ;
  \item \code{np.interp(...)} effectue l'interpolation linéaire de la section
        efficace sur les nombres d'onde demandés.
\end{itemize}

L'alias usuel \code{np} raccourcit les appels. Les opérations arithmétiques de
conversion s'appliquent aussi directement aux tableaux NumPy.

\subsection{\code{matplotlib.pyplot} sous l'alias \code{plt}}

Le sous-module \code{pyplot} fournit les commandes de tracé. Le script utilise
\code{plt.plot} pour la courbe, \code{plt.xlabel} et \code{plt.ylabel} pour les
axes, \code{plt.title} pour le titre, \code{plt.grid} pour la grille et
\code{plt.show} pour ouvrir la fenêtre graphique. Le reste du paquet
\code{matplotlib} n'est pas importé directement.

\subsection{Fonctions importées depuis \code{hapi}}

La forme \code{from hapi import ...} permet d'appeler directement trois
fonctions :

\begin{description}
  \item[\code{db\_begin(chemin)}] sélectionne le dossier utilisé par HAPI pour
        ses tables locales. \code{str(DOSSIER\_DONNEES)} convertit l'objet
        \code{Path} en chaîne de caractères avant de le transmettre à HAPI ;
  \item[\code{fetch(nom, molécule, isotopologue, min, max)}] interroge le
        serveur HITRAN et enregistre les raies dans une table locale ;
  \item[\code{absorptionCoefficient\_Voigt(...)}] calcule le spectre en
        sommant les raies, chacune étant représentée par un profil de Voigt.
\end{description}

Tous les imports Python du fichier sont effectivement utilisés ; aucun ne peut
être supprimé sans modifier le code.

\section{Dossier local de données HITRAN}

La constante suivante désigne le dossier situé à côté du script :

\begin{lstlisting}
DOSSIER_DONNEES = Path(__file__).resolve().parent / "donnees_hitran"
\end{lstlisting}

Après le téléchargement, HAPI y stocke notamment :

\begin{itemize}
  \item \code{CO2.data}, qui contient les raies spectroscopiques ;
  \item \code{CO2.header}, qui décrit la table locale.
\end{itemize}

\section{Fonctions du script}

\subsection{\code{convertir\_lambda\_m\_en\_nombre\_onde\_cm\_1}}

\begin{lstlisting}
def convertir_lambda_m_en_nombre_onde_cm_1(longueur_onde_m):
    return 1 / (longueur_onde_m * 100)
\end{lstlisting}

La fonction reçoit une longueur d'onde, ou un tableau de longueurs d'onde, en
mètres. Elle renvoie le nombre d'onde correspondant en
\si{\per\centi\metre}. Elle ne vérifie pas que les longueurs d'onde sont
strictement positives.

\subsection{\code{charger\_section\_co2}}

Cette fonction télécharge les raies, puis calcule la section efficace. Ses
paramètres par défaut sont :

\begin{center}
\begin{tabular}{lll}
\toprule
Paramètre & Valeur & Signification \\
\midrule
\code{nombre\_onde\_min} & \SI{550}{\per\centi\metre} & borne spectrale basse \\
\code{nombre\_onde\_max} & \SI{800}{\per\centi\metre} & borne spectrale haute \\
\code{pas\_nombre\_onde} & \SI{0.05}{\per\centi\metre} & pas de la grille HAPI \\
\code{temperature} & \SI{296}{\kelvin} & température du gaz \\
\code{pression\_atm} & \SI{1}{\atm} & pression du gaz \\
\bottomrule
\end{tabular}
\end{center}

L'appel de téléchargement est :

\begin{lstlisting}
fetch("CO2", 2, 1, nombre_onde_min, nombre_onde_max)
\end{lstlisting}

\code{"CO2"} est le nom de la table locale, \code{2} l'identifiant HITRAN du
\ce{CO2} et \code{1} l'identifiant de son isotopologue principal. Les
isotopologues minoritaires ne sont donc pas inclus dans ce prototype.

Le calcul spectral est ensuite réalisé par :

\begin{lstlisting}
nombres_onde, sections_cm2 = absorptionCoefficient_Voigt(
    SourceTables="CO2",
    WavenumberRange=[nombre_onde_min, nombre_onde_max],
    WavenumberStep=pas_nombre_onde,
    Environment={"T": temperature, "p": pression_atm},
    HITRAN_units=True,
)
\end{lstlisting}

\code{SourceTables} choisit la table locale, \code{WavenumberRange} et
\code{WavenumberStep} définissent la grille, et \code{Environment} fixe la
température en kelvins et la pression en atmosphères. Avec
\code{HITRAN\_units=True}, HAPI renvoie une section efficace en
\si{\centi\metre\squared\per\molecule}, et non un coefficient d'absorption en
\si{\per\centi\metre}.

La conversion vers le système international utilise
$\SI{1}{\centi\metre\squared}=10^{-4}\,\si{\metre\squared}$ :

\begin{lstlisting}
return nombres_onde, sections_cm2 * 1e-4
\end{lstlisting}

La fonction renvoie donc deux tableaux : les nombres d'onde en
\si{\per\centi\metre} et les sections efficaces en
\si{\metre\squared\per\molecule}.

\subsection{\code{calculer\_section\_co2}}

La fonction convertit d'abord les longueurs d'onde demandées en nombres
d'onde, puis effectue une interpolation linéaire :

\begin{lstlisting}
return np.interp(
    nombres_onde_demandes,
    nombres_onde,
    sections_co2,
    left=0,
    right=0,
)
\end{lstlisting}

\code{nombres\_onde} et \code{sections\_co2} sont les deux tableaux renvoyés
par \code{charger\_section\_co2}. Pour toute demande située hors de cette
grille, \code{left=0} et \code{right=0} imposent artificiellement une section
efficace nulle.

\subsection{\code{tracer\_section\_co2}}

Cette fonction coordonne le calcul et l'affichage. Elle crée 3000 longueurs
d'onde entre 12 et \SI{18}{\micro\metre}, interpole les sections efficaces et
convertit les longueurs d'onde de mètres en micromètres pour l'axe horizontal :

\begin{lstlisting}
plt.plot(longueurs_onde_m * 1e6, sections_interpolees)
\end{lstlisting}

Le facteur $10^6$ vient de
$\SI{1}{\metre}=10^6\,\si{\micro\metre}$. L'appel final à \code{plt.show()}
ouvre la fenêtre graphique.

\section{Exécution et limites à connaître}

Depuis le dossier du script :

\begin{lstlisting}[language=bash]
python 1_section_efficace_CO2_15um.py
\end{lstlisting}

La condition suivante empêche le tracé de s'exécuter lorsque le fichier est
seulement importé comme module :

\begin{lstlisting}
if __name__ == "__main__":
    tracer_section_co2()
\end{lstlisting}

Trois limites sont importantes pour interpréter le résultat :

\begin{itemize}
  \item la plage HITRAN de \SIrange{550}{800}{\per\centi\metre} correspond
        environ à \SIrange{12.50}{18.18}{\micro\metre}. Le graphique commence
        pourtant à \SI{12}{\micro\metre} : entre 12 et
        \SI{12.5}{\micro\metre}, la courbe est donc forcée à zéro par
        l'interpolation. Couvrir toute la figure demanderait une borne haute
        d'au moins \SI{834}{\per\centi\metre} ;
  \item seul l'isotopologue principal du \ce{CO2} est pris en compte ;
  \item les conditions par défaut, \SI{296}{\kelvin} et \SI{1}{\atm}, sont des
        conditions de référence et non un profil atmosphérique variant avec
        l'altitude.
\end{itemize}

En environnement sans interface graphique, \code{plt.show()} peut ne rien
afficher. Il faut alors employer un backend Matplotlib non interactif ou
remplacer cet appel par \code{plt.savefig(...)}.

\section{Place dans le modèle climatique}

Ce fichier sert à valider sur un seul gaz l'accès à HITRAN, le calcul du
profil spectral et les conversions d'unités. La section efficace peut ensuite
intervenir dans l'épaisseur optique d'une couche homogène :

\begin{equation}
  \tau(\lambda)=\sigma(\lambda)\,n\,\Delta z,
\end{equation}

où $n$ est la densité moléculaire en
\si{\molecule\per\metre\cubed} et $\Delta z$ l'épaisseur de la couche en
mètres. Dans un modèle sans diffusion et à l'équilibre thermodynamique local,
l'absorptivité spectrale de la couche s'écrit alors :

\begin{equation}
  \alpha(\lambda)=1-\exp[-\tau(\lambda)].
\end{equation}

D'après la loi de Kirchhoff, l'émissivité spectrale est égale à
l'absorptivité dans ces conditions. Le calcul étendu à plusieurs gaz est
réalisé dans \code{2\_sections\_efficaces\_gaz\_hitran.py}.

\end{document}
